\chapter{Testing, Deployment and Evaluation}
\label{ch:evaluation}

A system that claims strong security properties must be tested against them, not
merely described. This chapter evaluates \zkw on four axes: functional correctness,
established through an automated test suite; security, assessed against the threat
model of Chapter~\ref{ch:architecture}; design-goal attainment, judged against the
principles the system set for itself; and standing relative to prior art. It then
analyses the central claim of the thesis --- that distributing key custody across a
threshold network is a sounder foundation than concentrating it in trusted
hardware --- and closes with a deliberately frank account of the prototype's
limitations and the residual trust it cannot eliminate.

\section{Evaluation Methodology}
\label{sec:eval-method}

Three complementary methods are used. Functional correctness is established with a
deterministic unit-test suite that exercises every state-changing path in the
contracts, including adversarial and boundary cases, using Hardhat's ability to
manipulate block time and impersonate accounts. Security is evaluated analytically:
each threat enumerated in the threat model is revisited and traced to the
mechanism, or mechanisms, that mitigate it. Finally, the system is compared
qualitatively with the platforms and primitives surveyed in
Chapter~\ref{ch:related}, against the same criteria used there.

This mix is appropriate for a system of this kind. The on-chain logic is small,
self-contained, and deterministic, which makes it amenable to exhaustive unit
testing. The cryptographic guarantees, by contrast, rest on the security of
well-studied primitives (AES-256-GCM, secp256k1 ECDH, BLS threshold encryption,
zkTLS) and on the correct composition of those primitives, which is best argued
analytically rather than benchmarked. Where a property has not been fully
validated in the prototype, this is stated plainly rather than glossed.

\section{Functional Correctness: The Test Suite}
\label{sec:eval-tests}

The contracts are covered by 91 unit tests written with Hardhat, Chai, and
ethers, distributed as shown in Table~\ref{tab:test-coverage}. The suite covers
not only the happy paths but, more importantly, the access-control checks, the
revert conditions, the event emissions, the arithmetic of fee splitting, and the
time-dependent behaviour of the dead man's switch.

\begin{table}[h]
\centering
\caption{Automated test coverage by contract.}
\label{tab:test-coverage}
\begin{tabular}{p{5.2cm} c p{6.0cm}}
\toprule
\textbf{Contract} & \textbf{Tests} & \textbf{Principal concerns} \\
\midrule
\code{DeadMansSwitch}        & 28 & Heartbeat lifecycle, time-based triggering, deactivation, metadata updates \\
\code{WhistleblowerRegistry} & 33 & Three verification tiers, on-chain proof check, owner controls, anti-replay \\
\code{Marketplace}           & 30 & Listing, bidding, escrow, fee split, stealth payout, refunds, reentrancy guards \\
\midrule
\textbf{Total}               & \textbf{91} & \\
\bottomrule
\end{tabular}
\end{table}

The dead man's switch tests are the most subtle, because correctness there is a
function of time. Hardhat's \code{evm\_increaseTime} cheat lets the suite advance
the chain clock deterministically and assert that \isdeceased flips exactly at the
interval boundary, that a timely check-in resets a near-expired switch, and that a
deactivated switch never triggers even long after its nominal deadline
(Listing~\ref{lst:test-dms}). These are precisely the conditions on which the
entire release mechanism depends, so testing them directly is essential.

\begin{lstlisting}[language=TypeScript,caption={Time-dependent trigger behaviour (\code{test/DeadMansSwitch.ts}).},label={lst:test-dms}]
it("Should be true once the interval elapses", async function () {
  await increaseTime(ONE_DAY + 1);
  expect(await dms.isDeceased(owner.address)).to.be.true;
});

it("Should reset a near-expired switch so it is not deceased", async function () {
  await increaseTime(ONE_DAY - 10);
  await dms.connect(owner).checkIn();
  await increaseTime(10);
  expect(await dms.isDeceased(owner.address)).to.be.false;
});

it("Should be false after deactivation even if expired", async function () {
  await dms.connect(owner).deactivateSwitch();
  await increaseTime(ONE_DAY + 1);
  expect(await dms.isDeceased(owner.address)).to.be.false;
});
\end{lstlisting}

The registry tests are notable for exercising all three verification tiers and,
crucially, the security property that distinguishes the trustless path: that
verification accrues to the proof's \emph{owner}, not to whoever relays it. A
dedicated test relays a proof issued to one account from a different account and
asserts that only the owner becomes verified, confirming that a valid proof cannot
be replayed to credential an attacker's address (Listing~\ref{lst:test-registry}).
The suite also confirms that a rejected proof records nothing at all, that
duplicate identifiers are refused, and that owner-only functions revert for
non-owners with the expected custom error.

\begin{lstlisting}[language=TypeScript,caption={Verification accrues to the proof owner, not the relayer (\code{WhistleblowerRegistry.ts}).},label={lst:test-registry}]
it("Should accrue verification to the proof owner, not the relayer", async function () {
  // user1 relays a proof issued to user2.
  await registry.connect(user1).submitVerifiedProof(buildProof(IDENTIFIER, user2.address));
  expect(await registry.isVerified(user2.address)).to.be.true;
  expect(await registry.isVerified(user1.address)).to.be.false;
});

it("Should revert when the verifier rejects the proof", async function () {
  await mock.setShouldPass(false);
  await expect(registry.submitVerifiedProof(buildProof(IDENTIFIER, user1.address)))
    .to.be.revertedWith("MockReclaim: invalid proof");
  expect(await registry.isVerified(user1.address)).to.be.false;
  expect(await registry.proofHashExists(IDENTIFIER)).to.be.false;
});
\end{lstlisting}

The marketplace tests verify the financial logic end to end: that a bid below the
minimum is rejected, that a source cannot bid on their own listing, that accepting
a bid sends exactly the payout (bid minus the 2.5\% fee) to the supplied stealth
address and the fee to the fee recipient, that an accepted listing closes, and that
unaccepted bids can be fully refunded. The exact-balance assertions, which account
for gas in the withdrawal case, give confidence that no wei is lost or stranded and
that the fee arithmetic is correct. Together with the reentrancy guards verified
by construction, this gives the fund-handling code a solid correctness argument.

\section{Security Evaluation Against the Threat Model}
\label{sec:eval-security}

Table~\ref{tab:threat-mitigation} revisits the adversaries and threats of
Chapter~\ref{ch:architecture} and maps each to the mechanism that addresses it.
The recurring pattern is that no single component is trusted with a secret: an
adversary must defeat several independent systems --- the chain, the threshold
network, the storage layer --- simultaneously to compromise a vault.

\begin{table}[h]
\centering
\caption{Threats and their mitigations in \zkwns.}
\label{tab:threat-mitigation}
\begin{tabular}{p{4.6cm} p{8.8cm}}
\toprule
\textbf{Threat} & \textbf{Mitigation} \\
\midrule
Coerced or silenced source & Dead man's switch releases automatically on missed heartbeats; no living action required at release time. \\
Server seizure / takedown   & No server holds the data or the key; ciphertext on Arweave, key shares across the Lit network. \\
Storage operator reads data & Only AES-256-GCM ciphertext is stored; the operator never sees plaintext or the key. \\
Single custodian compromise & Threshold custody: no individual Lit node holds the whole key; an $n$-of-$m$ quorum is required to unseal. \\
Premature release            & Lit nodes independently evaluate \isdeceased on-chain; the key is withheld until the trigger is objectively true. \\
Forged credibility           & Verified badge is read from the registry; the trustless tier checks Reclaim witness signatures on-chain. \\
Proof replay                 & Verification accrues to the proof owner; duplicate identifiers are rejected. \\
Deanonymising payments       & One-time ERC-5564 stealth addresses break the payer--payee link on-chain. \\
Reentrancy on payouts        & Checks--effects--interactions plus \code{ReentrancyGuard} on all fund moves. \\
Secret exfiltration (Reclaim) & App secret confined to a Node-only route behind a provider allow-list; never sent to the browser. \\
\bottomrule
\end{tabular}
\end{table}

The strongest property to highlight is censorship resistance under the
``release-on-death'' condition. Because the trigger is public on-chain state and
the key is held by a decentralised quorum, there is no party --- not the platform
operator, not a hosting provider, not a government with jurisdiction over a single
server --- that can both prevent a legitimate release and is the sole route to one.
This is the qualitative leap over the centralised and even the partially
decentralised systems surveyed in Chapter~\ref{ch:related}.

\section{Attainment of Design Goals}
\label{sec:eval-goals}

Measured against the design principles set out in Chapter~\ref{ch:architecture},
the prototype substantially achieves what it set out to do, with one important
caveat on the release path. Table~\ref{tab:goal-attainment} summarises the
assessment.

\begin{table}[h]
\centering
\caption{Design-goal attainment.}
\label{tab:goal-attainment}
\begin{tabular}{p{0.9cm} p{6.6cm} p{5.8cm}}
\toprule
\textbf{ID} & \textbf{Principle} & \textbf{Status} \\
\midrule
P1 & Decentralised trust (no single point of failure) & Achieved: storage, custody, and trigger are independent decentralised systems. \\
P2 & No plaintext leaves the client                  & Achieved: AES-GCM in-browser; only ciphertext egresses. \\
P3 & No sensitive data on-chain                       & Achieved: contracts store hashes, timestamps, references only. \\
P4 & Mathematical rather than hardware trust          & Achieved (design + arming); release path implemented but not yet validated live. \\
P5 & Reward without deanonymisation                   & Achieved: one-time stealth addresses. \\
\bottomrule
\end{tabular}
\end{table}

The honest qualification concerns P4 at \emph{release} time. The arming path ---
generating a key, sealing it under the on-chain condition, storing only
ciphertext --- is implemented and exercised. The decryption path follows the
documented Lit~v8 \code{AuthManager} pattern but has not been run against a live
Naga deployment in this prototype, as the source comments record. The design is
complete and the property holds in principle; what remains is end-to-end
validation on a Lit-supported testnet, which is identified as immediate future
work in Chapter~\ref{ch:conclusions}.

\section{Comparison with Prior Art}
\label{sec:eval-comparison}

Returning to the comparison framework of Chapter~\ref{ch:related},
Table~\ref{tab:final-comparison} situates \zkw among representative systems. The
distinguishing feature is not any single capability --- others offer anonymous
submission, and others offer permanent storage --- but the \emph{combination} of an
automatic, censorship-resistant trigger, threshold key custody with no trusted
hardware, anonymous-yet-credible provenance, and unlinkable rewards, integrated in
one coherent design.

\begin{table}[h]
\centering
\caption{Qualitative comparison with representative systems.}
\label{tab:final-comparison}
\begin{tabular}{p{3.4cm} c c c c c}
\toprule
\textbf{System} & \textbf{Decentr.} & \textbf{Auto release} & \textbf{No-HW trust} & \textbf{Provenance} & \textbf{Anon. pay} \\
\midrule
SecureDrop            & No       & No  & n/a  & No  & No  \\
Classic DMS scripts   & Partial  & Yes & n/a  & No  & No  \\
TEE-based vaults       & Partial  & Yes & No   & No  & No  \\
\zkwns                & Yes      & Yes & Yes  & Yes & Yes \\
\bottomrule
\end{tabular}
\end{table}

The comparison should be read with care: the established systems are mature,
deployed, and battle-tested, whereas \zkw is a research prototype. The claim is not
that \zkw is ready to replace them today, but that its architecture occupies a
point in the design space --- maximal decentralisation of trust with no reliance on
trusted hardware --- that the others do not, and that this point is reachable with
primitives that exist now.

\section{Trust-Model Analysis: Threshold Cryptography versus Trusted Hardware}
\label{sec:eval-trust}

The intellectual core of this thesis is the argument that key custody for a
release-on-death system should rest on threshold cryptography rather than on a
Trusted Execution Environment. Having built both the surrounding system and the
threshold integration, the argument can now be made concretely.

A TEE concentrates trust in a single hardware vendor and a single physical chip.
Its confidentiality holds only so long as the enclave is not broken --- and the
literature documents a steady cadence of microarchitectural side-channel attacks
(Foreshadow, Plundervolt, SGAxe, and others) that have repeatedly extracted
secrets from production enclaves. A defender must trust the manufacturer's
implementation, its attestation service, and its microcode patch cadence. For a
system whose entire purpose is to protect a source from powerful adversaries, a
trust root that has been broken before, by exactly the class of well-resourced
attacker the system must withstand, is an uncomfortable foundation.

Threshold cryptography moves the trust root from a chip to a mathematical
assumption distributed over many independent operators. With an $n$-of-$m$ scheme,
an adversary must compromise a threshold number of nodes --- operated by different
parties, on different hardware, in different jurisdictions --- before learning
anything about the key. There is no single vendor to subvert and no single enclave
to attack. The trade-off is real and worth naming: the system now depends on the
liveness and honesty of a quorum of the network rather than on a chip, and it
inherits that network's availability and governance assumptions. But for the
specific property \zkw needs --- that no individual party can unilaterally read a
sealed secret --- distributing trust across many parties is a strictly stronger
position than concentrating it in one.

It is also why the implementation deliberately pins to Lit's MPC (Naga) path rather
than to newer TEE-backed runtimes. Adopting a hardware-backed Lit network would
re-import the very trust assumption the architecture set out to remove, collapsing
the distinction the thesis is built on. Keeping custody on the threshold path keeps
the design honest about where its trust ultimately lives.

\section{Cost and Performance Considerations}
\label{sec:eval-cost}

While not the focus of the evaluation, the prototype's economics are favourable for
its use case. Encrypted manifests under 100~KiB are stored on Arweave through
Turbo's free tier, so arming a vault incurs no storage cost and no funding
transaction --- a meaningful property for a source who may not hold tokens and cannot
risk a traceable funding step. On-chain costs are limited to the dead man's switch
operations: a one-time \code{createSwitch} and periodic \code{checkIn} calls, both
of which are simple storage writes. Encryption is a single AES-GCM pass executed
locally and is effectively instantaneous for the document-sized payloads the system
targets. The latency-sensitive operation is the Lit handshake for sealing and
unsealing the key, which involves a round trip to the threshold network; this is
acceptable for an arm-once, release-rarely workload, though it would matter for
higher-frequency applications.

The 100~KiB ceiling is a genuine constraint on payload size, discussed as a
limitation below; for larger leaks the natural extension is to store a large
ciphertext blob separately and keep only a compact, key-bearing manifest within the
free tier.

\section{Limitations}
\label{sec:eval-limitations}

A credible evaluation states what the system does not yet do. The principal
limitations of the current prototype are the following.

\textbf{The release path is not yet validated live.} As noted under P4, the Lit
decryption flow is implemented to the documented v8 pattern but has not been
exercised against a live Naga deployment. Until it is, the end-to-end
release-on-death guarantee is established by design and component testing rather
than by a full live run.

\textbf{Payload size is capped.} The free-tier storage path limits manifests to
100~KiB. Larger payloads require the split-blob extension sketched above, which is
designed but not implemented.

\textbf{Key custody depends on a specific network.} Lit evaluates access-control
conditions only on an allow-list of public chains, and the system's liveness is
tied to the Lit network and to Arweave remaining available. Decentralisation
reduces, but does not eliminate, dependence on these external networks; their own
governance and economics are outside the system's control.

\textbf{Metadata is not fully hidden.} Although no plaintext is on-chain, the
existence of a vault, its heartbeat interval, the timing of check-ins, and a
non-public recipient address are visible on-chain. A sophisticated observer could
draw inferences from this metadata --- for example, correlating a sudden cessation of
check-ins with external events. Stronger metadata privacy (e.g. via additional
mixing or private state) is future work.

\textbf{Liveness burden on the source.} The dead man's switch requires the source
to check in periodically; a missed check-in due to mere inconvenience, rather than
incapacitation, triggers a release. Tuning the interval trades off false triggers
against responsiveness, and the right setting is situational.

\textbf{In-memory key exposure.} The AES key is necessarily extractable in browser
memory during arming so that it can be sealed by Lit. This window is small and
local, but a fully compromised client device remains outside the threat model the
system can address.

These limitations are not fatal to the thesis; rather, they delineate the boundary
between what the prototype demonstrates --- that a no-trusted-hardware,
decentralised release-on-death architecture is buildable with today's primitives ---
and what a production system would additionally require. The next chapter
consolidates the contributions and lays out the path from prototype to that
production system.
